% ============================================
% Section: Related Work (MobiCom-oriented, network-protocol focused)
% ============================================

Co-LIC2 sits at the intersection of four communities: (i) wireless protocol
design for real-time streaming, (ii) edge-assisted collaborative SLAM,
(iii) multi-agent Gaussian and neural SLAM, and (iv) multi-sensor odometry.
We survey each and position our contributions.

\subsection{Wireless QoS and Multi-Stream Protocols}

The IEEE 802.11e amendment~\cite{ieee80211e} introduced EDCA (Enhanced
Distributed Channel Access) with four access categories (AC\_VO, AC\_VI, AC\_BE,
AC\_BK) differentiated by AIFSN, CWmin/max, and TXOP limits. Application-layer
QoS frameworks such as MPEG-DASH and WebRTC exploit EDCA for real-time video,
but \emph{no prior work extends EDCA-aware scheduling to the multi-modal traffic
mix of cooperative SLAM}: low-latency pose telemetry, bursty submap uploads,
keyframe references, and lightweight control messages. Co-LIC2 explicitly maps
these four traffic types to EDCA queues and adds application-layer FEC and
deadline-aware pacing.

Prior work on deadline-aware wireless scheduling (e.g., DRAGON~\cite{hou2023dragon}
in NSDI, PIAS~\cite{bai2015pias}) targets data-center or general-purpose
streaming; SwarmMap~\cite{xu2022swarmmap} and Map++~\cite{zhang2024mapp}
schedule SLAM tasks but do not differentiate traffic by latency criticality.
AdaMap~\cite{liu2023adamap} provides bounded end-to-end latency for cooperative
perception at the edge using priority queues, closest to our scheduling
philosophy, but targets LiDAR point-cloud fusion rather than Gaussian map
exchange. Co-LIC2's contribution is a protocol that simultaneously schedules
\emph{four traffic streams} with explicit per-stream QoS contracts and integrates
this scheduler with a Gaussian-specific compression optimizer.

\subsection{Edge-Assisted Collaborative SLAM}

The mobile computing community has led systems-level work on collaborative SLAM.
\textbf{SwarmMap}~\cite{xu2022swarmmap} (NSDI~2022) scales real-time
collaborative visual SLAM at the edge through change-log-based server--client
synchronization and priority-aware scheduling, serving $2\times$ more agents
than prior art. \textbf{Map++}~\cite{zhang2024mapp} (MobiCom~2024) introduces
user-participatory visual SLAM with metadata-based overlap assessment and
redundancy control, cutting traffic by $\approx$46\%. \textbf{Covins}~\cite{schmuck2021covins}
provides a decentralized-frontend, centralized-backend architecture for
VI-SLAM, and \textbf{Swarm-SLAM}~\cite{lajoie2024swarmslam} is a sparse fully
decentralized framework. These systems validate the architecture but target
feature-point maps over relatively benign links; they do not address the
bandwidth, compression, and streaming challenges of photo-realistic Gaussian
maps, nor do they design differentiated QoS for multi-modal traffic.

\subsection{Multi-Agent Gaussian and Neural SLAM}

Photo-realistic multi-agent mapping with 3D Gaussians is emerging rapidly.
\textbf{CGS-SLAM}~\cite{deambrogi2026cgsslam_mob} performs collaborative
Gaussian SLAM with metric monocular depth and VGGT-based alignment.
\textbf{GRAND-SLAM}~\cite{thomas2025grandslam_mob} achieves globally consistent
large-scale multi-agent Gaussian SLAM through local optimization and loop
closure. \textbf{MNE-SLAM}~\cite{deng2025mneslam} (CVPR~2025) provides
distributed neural SLAM with intra-to-inter loop closure and the INS multi-agent
dataset. \textbf{CP-SLAM}~\cite{hu2024cpslam} (NeurIPS~2024) is a collaborative
neural point-based SLAM. These works establish feasibility but target
RGB-D/monocular agents and do not model wireless constraints. Co-LIC2 is the
first to carry LiDAR+IMU per node (for metric scale and textureless-region
robustness) and to design the wireless protocol as a first-class system
component.

\subsection{Multi-Sensor Continuous-Time Odometry}

The per-node front-end builds on \textbf{Coco-LIC}~\cite{lang2023coco_mob}
(continuous-time B-spline LIV odometry with online time-offset calibration) and
\textbf{Gaussian-LIC2}~\cite{lang2026gaussian2_mob} (Coco-LIC + 3DGS for
real-time photo-realistic mapping). \textbf{FAST-LIVO2}~\cite{zheng2024fastlivo2_mob}
supports multi-cam/multi-LiDAR configurations in a discrete-keyframe framework.
Co-LIC2 generalizes Coco-LIC's residual stack to $N_c$ cameras and $N_l$ LiDARs
per node with online multi-offset estimation, which is essential for the
heterogeneous-node scenario and provides the continuous-time trajectory from
which submaps are carved at arbitrary temporal boundaries.

\subsection{Positioning}

In summary, prior work has solved the \emph{systems} problems of collaborative
SLAM for feature-point maps under real networks, and the \emph{vision} problems
of photo-realistic multi-agent Gaussian SLAM under benign links. Co-LIC2 is the
first system to design a \textbf{complete wireless protocol stack}---QoS-aware
multi-stream transport, deadline-driven compression scheduling, and hierarchical
degradation---for cooperative photo-realistic Gaussian SLAM. The protocol is
the paper's primary contribution to the MobiCom community.
